\section{Background}

% We review technologies at the intersection of developing optics devices and the fabrication of interactivity within HCI that inform our approach towards FDM printed optical devices. %. We explore prior methods to develop optical devices for sensing and dynamic display.
% , particularly in contexts such as textiles, interactive surfaces, and 3D printed interfaces. 
%We also highlight how prior research has enabled interactivity in passive 3D printed objects through optical encoding, material properties, and embedded geometries. These insights inform our approach to designing low-cost, multi-input/output systems using FDM. Finally, we examine various fabrication strategies for multi-component and modular objects, focusing on embedding optical and electronic components during or after printing. 

\subsection{Optical Devices in HCI}
Optical elements have been used extensively in Computational Fabrication for HCI research to enable sensing and information transmission. Examples include optical fiber touch sensors 
for interactive tables \cite{go2012object} and touch screens \cite{holz2013fiberio}, room-scale displays \cite{swaminathan2020optistructures}, textiles \cite{olwal2018braid, hashimoto2013lightcloth}, and grasp-sensitive surfaces \cite{wimmer2010flyeye}. To enhance flexibility, novel methods such as cellulose-based applications \cite{guridi2023light} and transparent elastomers \cite{yao2014integrating} have been explored. In the 3D printing domain, \textit{Printed Optics} \cite{printed_optics_2012} demonstrated the creation of optical channels through PolyJet printing for sensing and simple to advanced, curved displays \cite{brockmeyer2013papillon}, AR mapping \cite{fibar_2020}, and optimized routing \cite{clr_printed_optical_fibers_2014} of printed channels.
Although PolyJet printing provides significant advantages in creating arbitrarily shaped optical channels for assembly-free interactive device fabrication,
% easy component embedding, it is expensive and less accessible. In contrast, FDM is low-cost and accessible for desktop 3D printing, enabling makers and designers to design and develop their custom optical devices. 
% In this work, 
we demonstrate how FDM can be used to create optical channels for display and sensing. Furthermore, due to the fine-grained control over printing parameters and channel geometry, we show how light channeling can be varied through parameter changes. %, as shown by the material characterization.


\subsection{Optical Interactivity in Passive 3D Objects}
% Beyond creating optical pathways for light, 
% Researchers have explored methods for embedding interactivity into 3D printed objects—enabling input and output behaviors without relying on traditional electronics. 
% Prior work has incorporated photochromic \cite{jin2020photo} and thermochromic \cite{moon20243d} materials into 3D printed surfaces, allowing them to change color in response to light or heat. %These effects have also been combined in shape memory polymers to create context-aware visuals, such as QR codes that become visible only under specific conditions. Additionally, 
Interactivity is embedded through geometric or material encodings, only detectable by machines such as subsurface air cavities readable via structured light \cite{li2017aircode}, layered barcodes %using alternating material patterns 
\cite{maia2019layercode}, and infrared-fluorescent inlays for markerless tracking \cite{dogan2023brightmarker}. %Recent work has also used 
Lenticular structures have been used to produce view-dependent images %that shift displayed information based on the user’s perspective 
\cite{zeng2021lenticular}.
Surface anisotropy enables directional optical patterns (e.g., barcodes or icons) that are visible under specific lighting angles \cite{ma2023anisotag, Lu2025LumosX}.
% Internal channel design has enabled passive mechanical computation in 3D printed objects. Researchers have printed 
Embedded pipes act as light guides, or housings for post-processing elements \cite{series_of_tubes_2014}. %Air-driven logic gates \cite{airlogic_2022}, fabricated entirely from 3D printed geometries, can process user input and trigger physical or auditory outputs.
% Touch input has also been explored by embedding grid patterns beneath curved surfaces, allowing for high-resolution, multi-touch interactions directly on the object’s shell \cite{capacitive_touch_2024}.
While these prior works demonstrate approaches for embedding interactivity into 3D printed objects utilizing optics theory, most focus on isolated optical effects or tagging mechanisms. 
% Either the optical input and output are in one plane for almost parallel mapping\cite{printed_optics_2012}, or input to output mapping resolution is limited to 1:1 or 1:N \cite{series_of_tubes_2014}. In this work, we reframe optical channel routing as a design and interaction space, providing low cost ways of fabricating complex N:M input-output mapping.

In this work, we reframe optical channel routing as a design and interaction space, providing low cost ways of fabricating complex N:M input-output mappings.
Unlike previous methods where optical I/O is arranged in a single plane for nearly parallel mapping \cite{printed_optics_2012}, or limited to 1:1 or 1:N through auto-routing algorithms \cite{series_of_tubes_2014}, our approach uses the generated channels to segment the entire object into discrete and printable modules which allows users the flexibility for post-print modifications and reconfiguration. %To support FDM printing, channels must lie in-plane to create partitions.


\subsection{Fabricating Multi-Component Objects}
Printing optical channels for sensing and display within FDM 3D objects requires multi-part fabrication. 
Modular fabrication is an approach for enabling multi-component 3D printing,
% enhances design flexibility by 
decomposing a design into smaller 3D printed parts that are later assembled (e.g., \cite{series_of_tubes_2014}). 
In this way, integration of materials with various properties becomes relatively affordable,
compared to 
% technique enables designs that would be challenging  with 
conventional fabrication methods due to %. By dividing designs into smaller modules, designers can bypass 
printer size limitations and complexity constraints, making it possible to build large-scale or highly detailed objects through incremental assembly \cite{sun2021flextruss, huang20233d,moya2024single}. 
Modularity simplifies creating complex moving structures, separating functional and aesthetic sub-parts \cite{faruqi2023style2fab}, enabling interactive behaviors \cite{chen2021pneuseries}, %Additionally, modularity 
and accelerating prototyping by allowing parallel fabrication \cite{yildiz2024modular}, %when designs are divided into independent components, these parts can be produced simultaneously on multiple machines, 
significantly shortening overall build time \cite{nisser2021laserfactory, noodleprint2025}, while also producing structurally robust prints that topologically interlock the pieces \cite{voronoodle2022}. Modularity also allows multiple printers to print the object at the same time--- Voronoi-based schemes have been employed to divide and schedule printing work between different printers \cite{safezone2024, safezonestar2025}. Our work adopts modular FDM fabrication for optical devices to iteratively build complex objects, precisely place materials, and embed optical components during printing.
% , we utilize the modular design approach combined with hybrid assembly processes to ease the 


